\hypertarget{appendix-fdc7879e8bf2}{}\label{referencestart:appendix-fdc7879e8bf2}
\section{Quantization and compressed weights}\label{appendix-fdc7879e8bf2-quantization-and-compressed-weights}

\hypertarget{appendix-fdc7879e8bf2-quantization-and-compressed-weights}{}

\textbf{Curriculum depth:} Developed · \textbf{First seen:} C001 · \textbf{Latest development:} C006

\begin{allmightroledefinition}

\subsection{Definition and mechanism}\label{appendix-fdc7879e8bf2-definition-and-mechanism}

\hypertarget{appendix-fdc7879e8bf2-definition-and-mechanism}{}

Low-precision values and packed model weights reuse field-width, mask, shift, and bandwidth trade-offs.

\end{allmightroledefinition}

Quantization maps a larger numeric domain to a smaller discrete set, usually with information loss. Packing stores those discrete codes compactly; dequantization reconstructs approximations. Bit packing alone is lossless for valid codes, whereas quantization changes represented precision.

M006 develops affine scales and zero points, local group calibration, offset-coded INT4, one-bit row centroids, metadata overhead and error metrics. Packing preserves discrete codes; quantization deliberately loses information. See the worked five-value group in the continuous story.

\subsection{Prerequisites}\label{appendix-fdc7879e8bf2-prerequisites}

\hypertarget{appendix-fdc7879e8bf2-prerequisites}{}

\begin{itemize}
\tightlist
\item
  \href{../01-Bits-Programs-and-Reliable-Systems-P01/main.pdf\#appendix-1f9963e1a5d1}{Bit packing (Bits, Programs, and Reliable Systems, Part 1, Reference notes: Bit packing)} --- Packing places several constrained values into a carrier.
\item
  \href{../01-Bits-Programs-and-Reliable-Systems-P01/main.pdf\#appendix-2ed044358b82}{Arithmetic policies (Bits, Programs, and Reliable Systems, Part 1, Reference notes: Arithmetic policies)} --- A policy states the intended response to an unrepresentable answer: wrap modulo a base, report checked failure, widen storage, clamp to an endpoint, or grow arbitrary precision.
\end{itemize}

\subsection{Used by}\label{appendix-fdc7879e8bf2-used-by}

\hypertarget{appendix-fdc7879e8bf2-used-by}{}

\begin{itemize}
\tightlist
\item
  \href{../01-Bits-Programs-and-Reliable-Systems-P01/main.pdf\#appendix-4ccd7a775813}{One-bit centroid quantization (Bits, Programs, and Reliable Systems, Part 1, Reference notes: One-bit centroid quantization)} --- uses this mechanism to define its representation, error or validation contract.
\item
  \href{../01-Bits-Programs-and-Reliable-Systems-P01/main.pdf\#appendix-537171d69d27}{Quantization error metrics (Bits, Programs, and Reliable Systems, Part 1, Reference notes: Quantization error metrics)} --- uses this mechanism to define its representation, error or validation contract.
\end{itemize}

Connect later chunks here when they use this concept.

\subsection{Recall and next depth}\label{appendix-fdc7879e8bf2-recall-and-next-depth}

\hypertarget{appendix-fdc7879e8bf2-recall-and-next-depth}{}

Explain quantization and compressed weights without the label, work the example above, then state the assumption whose removal would change the result.

\textbf{Pending:} task-level accuracy and throughput experiments, alternative calibration objectives and target kernels; these were not measured in the teaching merge.

\subsection{Source and navigation}\label{appendix-fdc7879e8bf2-source-and-navigation}

\hypertarget{appendix-fdc7879e8bf2-source-and-navigation}{}

\href{../01-Bits-Programs-and-Reliable-Systems-P01/main.pdf\#01-foundations-i-ch001-choosing-what-information-to-keep}{Narrative development (Bits, Programs, and Reliable Systems, Part 1, Ch. 1)}

\href{../01-Bits-Programs-and-Reliable-Systems-P01/main.pdf\#01-foundations-i-ch001-from-a-signal-to-an-interpretation}{C001 teaching chapter (Bits, Programs, and Reliable Systems, Part 1, Ch. 1)} · \href{../01-Bits-Programs-and-Reliable-Systems-P03/main.pdf\#appendix-fd090d982a15}{Part definitions and mechanisms (Bits, Programs, and Reliable Systems, Part 3, Reference notes: Foundations)} · \href{../07-Vision-XR-and-Playable-Worlds-P01/main.pdf\#appendix-4ad47bdd269a}{Dependency Map (Vision, XR, and Playable Worlds, Part 1, Reference notes: Dependency Map)}
